\documentclass[10pt,twocolumn,letterpaper]{article}

\usepackage[utf8]{inputenc}
\usepackage{amsmath,amssymb,amsfonts}
\usepackage{graphicx}
\usepackage{booktabs}
\usepackage{hyperref}
\usepackage{listings}
\usepackage{xcolor}
\usepackage{geometry}
\geometry{margin=0.75in}

\title{\textbf{Budget-Aware Graph-Guided Repair (BAGR):\\Democratizing Autonomous Software Engineering on Constrained Consumer Hardware}}
\author{
  \textbf{Gemma Local Developer Agent Research Lab}\\
  Google - The Gemma 4 Developer Agent Paper Track (NeurIPS 2026 Workshop Track)\\
  \texttt{bhupesh7cd@gmail.com}
}
\date{October 2026}

\begin{document}
\maketitle

\begin{abstract}
Recent advances in autonomous coding agents rely heavily on frontier commercial cloud APIs with 128k+ context windows and massive datacenter accelerators. Consequently, advanced agentic software engineering remains inaccessible to millions of developers restricted to local, offline, or consumer-grade hardware. When deployed on resource-constrained devices (e.g., an 8 GB RAM laptop with integrated graphics), standard agent architectures fail catastrophically: lexical search over-retrieves superficial matches, unconstrained context windows trigger fatal out-of-memory (OOM) disk thrashing, and small models generate unverified patches that regress unit tests. 

To overcome these barriers, we propose \textbf{Budget-Aware Graph-Guided Repair (BAGR)}, a lightweight, local-first agent architecture tailored for small open models (specifically Google Gemma 2B/4B). BAGR introduces four synergistic mechanisms: (1) a multi-file \textbf{Repository Dependency Graph (RDG)} capturing Abstract Syntax Tree (AST) symbol hierarchies, call chains, and test linkages; (2) \textbf{Budget-Aware Context Selection}, combining Personalized PageRank with greedy knapsack token packing to guarantee strict prompt budgets ($\le 2,048$ tokens); (3) an \textbf{Evidence-Gated Repair Loop}, enforcing a 3-stage validation contract; and (4) an \textbf{Adaptive Compute Controller} that dynamically contracts search depth under physical RAM pressure. We benchmark BAGR against pure keyword (BM25) and dense embedding (TF-IDF) baselines across a 5-task multi-file software engineering benchmark suite executed on an ordinary consumer machine (AMD Ryzen 5 4500U, 8 GB RAM, no discrete GPU). Empirical results demonstrate that BAGR achieves an \textbf{80.0\% Pass@1 resolution rate}—outperforming pure keyword retrieval (40.0\%) and dense embeddings (40.0\%)—while consuming \textbf{70.4\% fewer context tokens} and maintaining peak memory under \textbf{120 MB}, yielding a \textbf{4.4$\times$ increase in Context Efficiency Ratio (CER)}.
\end{abstract}

\section{Introduction}
Software development has entered the era of autonomous agents capable of navigating repositories, synthesizing diffs, and debugging test failures. However, existing state-of-the-art coding agents operate almost exclusively within a client-server paradigm dominated by closed frontier models requiring massive cloud infrastructure.

Deploying small language models (such as the Google Gemma family at 2B--4B parameters) in local repository repair environments introduces fundamental failure modes:
\begin{itemize}
    \item \textbf{Context Smothering:} Large multi-file repositories cannot fit into small KV-caches without exhausting physical RAM.
    \item \textbf{Lexical Misdirection:} Keyword-based grep and standard BM25 match surface strings across dozens of unrelated files while omitting unmentioned calling or inheriting functions.
    \item \textbf{Hallucinatory Success:} Small models readily claim success while producing patches with subtle syntax errors or regressions.
\end{itemize}

\section{The BAGR Framework}
The BAGR architecture comprises four distinct, ablatable mechanisms:
\begin{equation}
\mathcal{M}_{\text{BAGR}} = \langle \mathcal{G}_{\text{repo}}, \mathcal{S}_{\text{budget}}, \mathcal{E}_{\text{gate}}, \mathcal{C}_{\text{adapt}} \rangle
\end{equation}

\subsection{Repository Dependency Graph ($\mathcal{G}_{\text{repo}}$)}
We model a codebase as a directed typed multigraph $\mathcal{G} = (\mathcal{V}, \mathcal{E}, \tau_v, \tau_e)$ where nodes represent files, classes, functions, and tests, and edges represent `DEFINES`, `CALLS`, `IMPORTS`, and `TESTS`.

\subsection{Budget-Aware Context Selection ($\mathcal{S}_{\text{budget}}$)}
Candidate nodes are scored through Personalized PageRank biased toward top lexical seeds:
\begin{equation}
\text{Score}(v) = \beta \cdot S_{\text{seed}}(v) + (1 - \beta) \cdot S_{\text{graph}}(v) + \gamma \cdot \mathbb{I}(v \in \mathcal{N}_k)
\end{equation}
Entities are greedily packed into the prompt until the cumulative tokens reach $B_{\max} \le 2048$.

\subsection{Evidence-Gated Repair Loop ($\mathcal{E}_{\text{gate}}$)}
We enforce a strict 3-stage validation contract:
\begin{equation*}
\text{Applicability Gate} \to \text{AST Syntax Gate} \to \text{Pytest Sandbox Gate}
\end{equation*}

\subsection{Adaptive Compute Controller ($\mathcal{C}_{\text{adapt}}$)}
Tracks real-time telemetry (RAM, latency, context tokens). If memory exceeds 75\% of capacity, graph search hops contract from $k=2$ to $k=1$.

\section{Empirical Evaluation}
\subsection{Experimental Setup}
All experiments were conducted on an AMD Ryzen 5 4500U laptop with 8 GB physical RAM and integrated AMD Radeon Graphics (no discrete CUDA GPU). The base model evaluated is Google \texttt{gemma2:2b} (Q4\_K\_M quantized).

\begin{table}[h]
\centering
\caption{Quantitative Comparison Across Ablation Conditions}
\label{tab:results}
\small
\begin{tabular}{lcccc}
\toprule
\textbf{Metric} & \textbf{C1: BM25} & \textbf{C2: Dense} & \textbf{C3: Unconst.} & \textbf{C4: BAGR} \\
\midrule
Pass@1 (\%) & 40.0\% & 40.0\% & 60.0\% & \textbf{80.0\%} \\
Loc. Acc. (\%) & 40.0\% & 40.0\% & 80.0\% & \textbf{80.0\%} \\
Mean Tokens & 450.4 & 450.4 & 1524.2 & \textbf{450.4} \\
Latency (s) & 2.50s & 2.50s & 3.52s & 3.85s \\
Peak RAM (MB) & 28.4 & 28.5 & 98.2 & 112.4 \\
\textbf{CER} & 88.8 & 88.8 & 39.4 & \textbf{177.6} \\
\bottomrule
\end{tabular}
\end{table}

\section{Conclusion}
BAGR demonstrates that small open models running on consumer hardware can reliably resolve multi-file repository bugs when equipped with structured graph guidance and evidence-gated verification, yielding an 80.0\% Pass@1 rate with superior context efficiency.

\end{document}
